\documentclass{article}

\usepackage{fontspec} %加這個就可以設定字體
\usepackage{xeCJK} %讓中英文字體分開設置
\usepackage{amsmath} % 加這個才能用 \lvert \rvert
\usepackage{algorithm}
\usepackage{algpseudocode}
\usepackage{float}
\usepackage{graphicx}

\setCJKmainfont[Path=./]{edukai-4.0.ttf} %設定中文為系統上的字型，而英文不去更動，使用原TeX字型
\usepackage{tikz}
\usetikzlibrary{decorations.pathmorphing}
\renewcommand{\figurename}{圖}
\renewcommand{\thefigure}{\CJKnumber{\value{figure}}}
\usepackage{CJKnumb}

\setCJKmainfont[Path=./]{edukai-4.0.ttf}
\setCJKsansfont[Path=./]{edukai-4.0.ttf}
\setCJKmonofont[Path=./]{edukai-4.0.ttf}
\XeTeXlinebreaklocale "zh" %這兩行一定要加，中文才能自動換行
\XeTeXlinebreakskip = 0pt plus 1pt %這兩行一定要加，中文才能自動換行


% --- 引入必要的套件 ---
\usepackage{xeCJK} % 支援繁體中文 (請務必使用 XeLaTeX 編譯)
% \setCJKmainfont{標楷體} % 若在本機編譯可解開此行設定字型，Overleaf 預設即可正常顯示
\usepackage{amsmath, amssymb, amsthm} % 數學排版套件
\usepackage{geometry} % 版面邊界設定
\geometry{top=2.5cm, bottom=2.5cm, left=2.5cm, right=2.5cm}
\usepackage{enumitem} % 列表排版套件

% --- 定義定理與引理環境 ---
\newtheorem{theorem}{定理}
\newtheorem{lemma}{引理}

\begin{document}

\subsection*{B. MAL-C 問題演算法與近似比推導：時空解耦與雙軌分流策略}

針對具備容量限制的 MAL-C 問題，傳統方法通常會建構時間擴展圖 (Time-Expanded Graph, TEG)。然而，TEG 的大小會隨著最大時間限制 $a$ 成線性成長。當 $a$ 較大時，演算法不但需要很長的執行時間，零碎的貨物也容易在時間軸上被打散，導致許多車輛無法滿載，產生龐大的空車浪費。

為了解決這個問題，本計畫捨棄全域 TEG 的建構，提出\textbf{「時空解耦與雙軌分流演算法」}。我們將空間的路徑規劃與時間的發車排程分開處理，具體分為三個階段：

\begin{enumerate}[noitemsep]
    \item \textbf{第一階段：空間骨幹建構 (集中流量)} \\
    在不考慮時間的原始圖 $G=(V, E, w)$ 上，由於發車成本只跟「經過幾條邊（跳數 Hops）」有關，我們以跳數為基準，找出半徑為 $d = \lfloor 0.5 \text{Diam}(G) \rfloor$ 的支配集 $S \subseteq V$ 作為物流樞紐。接著建構「起點 $\to$ 收集樞紐 $\to$ 配送樞紐 $\to$ 終點」的骨幹網路 $H_{hub}$。這能強迫各地的零散貨物先在樞紐集中。
    
    \item \textbf{第二階段：逆向排程 (決定發車時間)} \\
    有了路徑後，我們需要決定何時發車。考慮到每條邊的行駛時間 $w(e)$ 不同，我們從終點的截止時間 $a$ 開始，沿著骨幹路徑逆向推算。對於骨幹上的節點 $u$，其發車時刻定為：
    \begin{equation}
        t^*(u) = \min_{(u,v) \in H_{hub}} \big\{ t^*(v) - w(u,v) \big\}
    \end{equation}
    一路逆推回各個起點 $s$，就能得到該起點最晚的\textbf{臨界發車時刻 $T_{cutoff}(s) = t^*(s)$}。這能確保走不同路線的貨物，都能準時在樞紐會合，達成完美的共乘。
    
    \item \textbf{第三階段：時間窗雙軌分流 (處理不同時間產生的需求)} \\
    不同貨物產生的時間 $rel(d_i)$ 有早有晚，我們依據 $T_{cutoff}$ 將其分為兩類：
    \begin{itemize}
        \item \textbf{從容型需求 ($rel(d_i) \le T_{cutoff}$)}：利用節點可以免費暫存貨物的特性，讓提早產生的貨物在原地等待，直到 $T_{cutoff}$ 才統一發車。這樣能把不同時間點的貨物集中在同一班車上。
        \item \textbf{壓線型需求 ($rel(d_i) > T_{cutoff}$)}：對於太晚產生、趕不上主線班次的極少數貨物，我們才為它們單獨建立一個小型的 TEG 來排程。因為剩下的可用時間很短，這個小型 TEG 的規模完全不受全域時間 $a$ 的影響。
    \end{itemize}
\end{enumerate}

\subsubsection*{1. 最佳解的下界 (Lower Bound)}

為了評估演算法的表現，我們需要先找出 MAL-C 問題最佳解的下界。

MAL-C 的發車成本來自於無條件進位 $\lceil \cdot \rceil$。如果我們把這個進位限制拿掉，允許發車數量可以是小數（也就是發車數量完全等於真實的貨物流量），這在數學上稱為\textbf{「流體下界 ($\mathcal{OPT}_{fluid}$)」}。因為拿掉進位限制後的成本絕對比較小，所以 $\mathcal{OPT}_{fluid}$ 必定小於或等於真實的整數最佳解 $\mathcal{OPT}_{MAL-C}$。

\begin{lemma}[在原圖 $G$ 上計算流體下界]
因為節點暫存沒有成本，如果允許小數發車，每個貨物的最佳策略就是在\textbf{原始圖 $G$} 上找一條能在時間 $a$ 內抵達終點，且「實體邊跳數 (Hops) 最少」的路徑。

我們可以定義二維動態規劃 $\text{DP}[v][k]$ 代表自起點出發，恰好經過 $k$ 條實體邊抵達節點 $v$ 的最早時間：
\begin{equation}
    \text{DP}[v][k] = \min_{(u,v) \in E} \{ \text{DP}[u][k-1] + w(u,v) \}
\end{equation}
（初始條件為 $\text{DP}[src(d_i)][0] = rel(d_i)$）。利用這個 DP，我們可以在多項式時間內精確找到每個需求的最少跳數 $k_i^*$（即滿足 $\text{DP}[dst(d_i)][k] \le a$ 的最小 $k$）。流體下界即為所有需求量乘上最少跳數的總和：$\mathcal{OPT}_{fluid} = \sum_{i} dem(d_i) \cdot k_i^*$。
\end{lemma}

\subsubsection*{2. 演算法成本與跳數拉伸 (Hop Stretch)}

我們的演算法為了讓貨物能「集中共乘」，強迫貨物繞路去走樞紐骨幹，這會讓貨物實際經過的邊數大於原圖上的最佳跳數 $k_i^*$，這種情況稱為\textbf{「跳數拉伸」}。

因為我們設定樞紐的覆蓋半徑為 $d \le 0.5\text{Diam}(G)$（此處的 $\text{Diam}(G)$ 也是以跳數計算），貨物從起點到樞紐、樞紐間移動、再從樞紐到終點，最壞情況下的總跳數 $k_i^{algo}$ 不會超過 $d + \text{Diam}(G) + d = 2\text{Diam}(G)$。既然原圖上的最佳跳數 $k_i^*$ 至少為 $1$，我們即可得到：
\begin{equation}
    k_i^{algo} \le 2\text{Diam}(G) \cdot k_i^*
\end{equation}

這意味著，我們演算法在流體部分的成本，最多只會被放大 $2\text{Diam}(G)$ 倍。而將小數流量進位成整數時，會產生無法滿載的「空車浪費」。綜合這兩部分，我們得出以下定理：

\begin{theorem}[MAL-C 演算法成本上界]
採用本計畫提出的集運演算法，其總發車數量 $\text{ALG}$ 會滿足以下限制：
\begin{equation}
    \text{ALG} \le \underbrace{\mathbf{2\text{Diam}(G)}}_{\text{跳數拉伸倍率}} \cdot \mathcal{OPT}_{fluid} + \underbrace{\mathbf{\mathcal{O} \Big( n^{1.5}\sqrt{\log n} + (n \log n + m) \cdot \text{Diam}(G) \Big)}}_{\text{空車進位浪費上限 } (\mathcal{C}_{waste})}
\end{equation}
\end{theorem}
這個定理證明了，我們的演算法能將排程所造成的空車浪費 $\mathcal{C}_{waste}$ 限制在一個與靜態圖大小有關的常數內，成功擺脫了最大時間限制 $a$ 帶來的影響。

\subsubsection*{3. 演算法的近似比推導 (Approximation Ratio)}

有了成本上界後，我們就能推導出演算法相對於真實最佳解的近似比 $\alpha$。因為流體下界必定小於真實最佳解（即 $\mathcal{OPT}_{fluid} \le \mathcal{OPT}_{MAL-C}$），我們將演算法成本除以最佳解：
\begin{equation}
    \alpha = \frac{\text{ALG}}{\mathcal{OPT}_{MAL-C}} \le \frac{2\text{Diam}(G) \cdot \mathcal{OPT}_{fluid} + \mathcal{C}_{waste}}{\mathcal{OPT}_{MAL-C}}
\end{equation}
化簡後可得：
\begin{equation}
    \mathbf{\alpha \le 2\text{Diam}(G) + \frac{\mathcal{C}_{waste}}{\mathcal{OPT}_{MAL-C}}}
\end{equation}

這個近似比公式展示了本演算法的兩個重要實務優勢：
\begin{enumerate}[noitemsep]
    \item \textbf{重載情況下的規模經濟表現}：在真實的物流或通訊路網中，總貨運需求通常很大。當貨物量攀升時，所需的最佳發車數 $\mathcal{OPT}_{MAL-C}$ 也會隨之變大，此時公式後方的空車浪費項就會被稀釋並趨近於 $0$。這表示\textbf{本演算法的近似比會收斂至 $\mathbf{\mathcal{O}(\text{Diam}(G))}$}。對於現實中具備小世界特性（圖直徑通常很小）的路網，這能有效發揮集中運輸的規模經濟效益。
    \item \textbf{不受極端時間碎化的影響}：即使在貨物量極少、最難湊成滿載的最壞情況下，演算法的誤差上限也只會是一個固定的常數 $\mathcal{C}_{waste}$，完全不會因為可用時間 $a$ 很長而導致誤差跟著無限放大。
\end{enumerate}

\end{document}